%%%%%%%%%%%%%%%%%%%%%%%%%%%%%%%%%%%%%%%%%
% Professional Formal Letter
% LaTeX Template
% Version 2.0 (12/2/17)
%
% This template originates from:
% http://www.LaTeXTemplates.com
%
% Authors:
% Brian Moses
% Vel (vel@LaTeXTemplates.com)
%
% License:
% CC BY-NC-SA 3.0 (http://creativecommons.org/licenses/by-nc-sa/3.0/)
%
%%%%%%%%%%%%%%%%%%%%%%%%%%%%%%%%%%%%%%%%%

%----------------------------------------------------------------------------------------
%	PACKAGES AND OTHER DOCUMENT CONFIGURATIONS
%----------------------------------------------------------------------------------------

\documentclass[10pt,letterpaper]{letter} % Set the font size (10pt, 11pt and 12pt) and paper size (letterpaper, a4paper, etc)

%---------------------------------------------------------
%   PACKAGES AND OTHER DOCUMENT CONFIGURATIONS
%---------------------------------------------------------
\usepackage{graphicx}   % Required for including pictures
\usepackage[T1]{fontenc} % Output font encoding for international characters
\usepackage[utf8]{inputenc} % Required for inputting international characters
\usepackage{gfsdidot}   % Use the GFS Didot font: http://www.tug.dk/FontCatalogue/gfsdidot/
\usepackage{microtype}  % Improves typography
\pagestyle{empty}       % Suppress headers and footers
\usepackage{xcolor}

\setlength\parindent{1cm} % Paragraph indentation

% Create a new command for the horizontal rule in the document which allows thickness specification
\makeatletter
\newcommand{\vhrulefill}[1]{\leavevmode\leaders\hrule\@height#1\hfill \kern\z@}
\makeatother

%---------------------------------------------------------
%   DOCUMENT MARGINS
%---------------------------------------------------------
\usepackage{geometry} % Required for adjusting page dimensions

\geometry{
    top=1cm, % Top margin
    bottom=1cm, % Bottom margin
    left=1in, % Left margin
    right=1in, % Right margin
    %showframe, % Uncomment to show how the type block is set on the page
}

%---------------------------------------------------------
%   DEFINE CUSTOM COMMANDS
%---------------------------------------------------------
\newcommand{\logo}[1]{\renewcommand{\logo}{#1}}
\newcommand{\Who}[1]{\renewcommand{\Who}{#1}}
\newcommand{\Title}[1]{\renewcommand{\Title}{#1}}
\newcommand{\headerlineone}[1]{\renewcommand{\headerlineone}{#1}}
\newcommand{\headerlinetwo}[1]{\renewcommand{\headerlinetwo}{#1}}
\newcommand{\authordetails}[1]{\renewcommand{\authordetails}{#1}}

%---------------------------------------------------------
%   AUTHOR DETAILS STRUCTURE
%---------------------------------------------------------
\newcommand{\authordetailsblock}{
    \begin{flushright}
        % \footnotesize % Use a smaller font size for the details
        % \Who \\ % Author name
        % \textit{\authordetails} % The author details text, all italicised
        \authordetails % The author details text, all italicised
    \end{flushright}
}

%---------------------------------------------------------
%   HEADER STRUCTURE
%---------------------------------------------------------
\address{
    \includegraphics[width=9cm]{\logo} % Include the logo of author institution
    \hspace{0.0\textwidth} % Position of the institution logo, increase to move left, decrease to move right
    \vskip -0.09\textheight~\\ % Position of the large header text in relation to the institution logo, increase to move down, decrease to move up
    \Large\hspace{0.2\textwidth}\headerlineone\hfill ~\\[0.006\textheight] % First line of institution name, adjust hspace if your logo is wide
    \hspace{0.2\textwidth}\headerlinetwo\hfill \normalsize % Second line of institution name, adjust hspace if your logo is wide
    \makebox[0ex][r]{\textbf{School of Physics and Astronomy}} \hspace{-0.018\textwidth} % Print author name and title with a little whitespace to the right
    ~\\[-0.005\textheight] % Reduce the whitespace above the horizontal rule
    \hspace{0.072\textwidth}{\color{lightgray}\vhrulefill{1pt}} \\[-0.01\textheight] % Horizontal rule, adjust hspace if your logo is wide and \vhrulefill for the thickness of the rule
    \authordetailsblock % Include the letter author's details on the right side of the page under the horizontal rule
    \hspace{-0.28\textwidth} % Horizontal position of the author details block, increase to move left, decrease to move right
    \vspace{-0.05\textheight} % Move the date and letter content up for a more compact look
}

%---------------------------------------------------------
%   COMPOSE THE ENTIRE HEADER
%---------------------------------------------------------
\renewcommand{\opening}[1]{
    {\raggedright\fromaddress\vspace{0.05\textheight} \\ % Print the header and from address here, add whitespace to move date down
    % \hspace*{\longindentation}\today\hspace*{\fill}\par} % Print today's date, remove \today to not display it
    \hspace*{\longindentation}June 26, 2017\hspace*{\fill}\par} % Print today's date, remove \today to not display it
    {\raggedright \toname \\ \toaddress \par} % Print the to name and address
    \vspace{1cm} % White space after the to address
    \noindent #1 % Print the opening line
}

%---------------------------------------------------------
%   SIGNATURE STRUCTURE
%---------------------------------------------------------
% \signature{\Who\Title} % The signature is a combination of the author's name and title 
\signature{R. P. Anderson, for the authors} % The signature is a combination of the author's name and title

\renewcommand{\closing}[1]{
    \vspace{2.5mm} % Some whitespace after the letter content and before the signature
    \noindent % Stop paragraph indentation
    \hspace*{\longindentation} % Move the signature right to the value of \longindentation
    \parbox{\indentedwidth}{
        \raggedright
        #1 \\% Print the signature text
        \includegraphics[width=4cm]{rpa_signature.png} \\
        % \vskip 1.65cm % Whitespace between the closing text and author's name for a physical signature
        \fromsig % Prints the value of \signature{}, i.e. author name and title
    }
}
 % Include the file that specifies the document structure

\longindentation=0pt % Un-commenting this line will push the closing "Sincerely," and date to the left of the page

%----------------------------------------------------------------------------------------
%	YOUR INFORMATION
%----------------------------------------------------------------------------------------

\Who{Dr Russell Anderson} % Your name

\Title{} % Your title, leave blank for no title

\authordetails{
	% School of Physics and Astronomy\\ % Your department/institution
	19 Rainforest Walk\\ % Your address
	Monash University \\
	Victoria, Australia 3800\\ % Your city, zip code, country, etc
	% russell.anderson@monash.edu\\ % Your email address
	% (+61) 3-9905-5943 \\ % Your phone number
	% URL: physics.monash.edu % Your URL
}

%----------------------------------------------------------------------------------------
%	HEADER CONTENTS
%----------------------------------------------------------------------------------------

\logo{monash_logo.pdf} % Logo filename, your logo should have square dimensions (i.e. roughly the same width and height), if it does not, you will need to adjust spacing within the HEADER STRUCTURE block in structure.tex (read the comments carefully!)

\headerlineone{} % Top header line, leave blank if you only want the bottom line

\headerlinetwo{} % Bottom header line

%----------------------------------------------------------------------------------------

\begin{document}

%----------------------------------------------------------------------------------------
%	TO ADDRESS
%----------------------------------------------------------------------------------------

\begin{letter}{
	Physical Review Letters\\
	Editorial Office\\
	1 Research Road \\
	Ridge, NY 11961-2701\\
}

%----------------------------------------------------------------------------------------
%	LETTER CONTENT
%----------------------------------------------------------------------------------------

\opening{Dear Editors,}

We are pleased to submit our manuscript 'Continuously observing a dynamically decoupled spin-1 quantum gas' to \textit{Physical Review Letters}.
In it we observe in real-time the continuous dynamical decoupling of a spin-1 quantum gas, probing its spectrum, coherences, and coupling strengths.
Dynamical decoupling a system from its environment is not only applicable to quantum information but has broad application to quantum metrology and quantum state estimation.
Yet few systems permit dynamical decoupling to be observed in real-time; weak measurements of superconducting qubits and nitrogen-vacancy (NV) centers require thousands of repeated preparation and detection cycles to measure what we do in a single-preparation of the quantum gas.

This work weds clock-state metrology with continuous measurement, ingredients for new experiments that emulate quantum magnetism in solid-state systems, and advancing a platform for measuring oscillating magnetic fields in noisy environments.
More broadly still, the decoupled microscale quantum gas offers magnetic sensitivity in a tunable kilohertz band, persistent over many milliseconds: this is the length scale, frequency, and duration relevant to sensing biomagnetic phenomena such as neural spike trains.

We advance spectrogram analysis as a time-frequency reduction of the measurement record to extract rich information about the quantum state and Hamiltonian simultaneously.
Spectrograms have been applied widely, from visualizing gravitational wave signals to analyzing human speech.
We `listen' to the precessing spins of the decoupled quantum gas, and use spectrogram analysis to parse the recorded tones as a signature of dynamical decoupling.

Our manuscript forms a co-submission with D. Trypogeorgos \textit{et al.}, who present a complementary characterization of the high-order decoupling manifest in this system as insensitivity of the states to magnetic field variations.
A distinguishing focus of our work is the continuous weak measurement, which could give rise to new forms of quantum sensing exploting synchronous detection and feedback.
Additionally, we demonstrate the ability to probe dressed-state coupling strengths without driving transitions, a form of Hamiltonian parameter estimation.
Our work is further distinguished by new analytic results identifying where the optimum decoupling occurs, applicable to any multi-level system.

\closing{Sincerely,}

%----------------------------------------------------------------------------------------
%	OPTIONAL FOOTNOTE
%----------------------------------------------------------------------------------------

% Uncomment the 4 lines below to print a footnote with custom text
%\def\thefootnote{}
%\def\footnoterule{\hrule}
%\footnotetext{\hspace*{\fill}{\footnotesize\em Footnote text}}
%\def\thefootnote{\arabic{footnote}}

%----------------------------------------------------------------------------------------

\end{letter}

\end{document}
